\chapter{LaTeX Docker}
\ChapterLabel{LatexDocker}
\index{LaTeX}
\index{Docker}
\index{Docker!TeX Live}
\index{Overleaf}

\section{Overview}
Overleaf compiles LaTeX projects on the server using a TeX Live installation.
This chapter reproduces that idea locally: a Docker container that bundles TeX
Live and compiles a simple LaTeX document into a PDF, so the build is the same
on every machine and requires no local TeX installation.

\subsection{Team Contributions}
\begin{tabularx}{\linewidth}{L{4.0cm}Y}
\toprule \textbf{Member} & \textbf{Contribution}\\ \midrule
Rifat Rahman Khan (RRK) & Wrote this chapter: designed the TeX Live Dockerfile, the sample LaTeX document, and the compile helper script.\\ \bottomrule
\end{tabularx}

\section{Project Layout}
\begin{minted}[fontsize=\small,frame=lines]{text}
latex-docker/
|-- Dockerfile
|-- hello.tex
|-- compile.sh
`-- output/          (created at run time; holds hello.pdf)
\end{minted}

\section{The LaTeX Source}
\begin{listing}[H]
\inputminted[fontsize=\small,breaklines,frame=lines,linenos]{latex}{latex-docker/hello.tex}
\caption{\texttt{hello.tex}: the document compiled by the container.}
\end{listing}

\section{The Dockerfile}
The image is based on Debian slim and installs a minimal TeX Live selection
(\texttt{texlive-latex-base}, \texttt{texlive-latex-recommended},
\texttt{texlive-latex-extra}, \texttt{texlive-fonts-recommended}) plus
\texttt{latexmk}, the same build driver Overleaf uses. The full
\texttt{texlive-full} package is much larger (several GB) and is unnecessary for
a simple document; swap it in if the document needs additional packages.

\begin{listing}[H]
\inputminted[fontsize=\small,breaklines,frame=lines,linenos]{docker}{latex-docker/Dockerfile}
\caption{\texttt{Dockerfile} for the TeX Live compile container.}
\end{listing}

Key design points:
\begin{itemize}
  \item \texttt{DEBIAN\_FRONTEND=noninteractive} prevents \texttt{apt} from
        prompting during the build.
  \item \texttt{--no-install-recommends} and removing \texttt{/var/lib/apt/lists}
        keep the image small.
  \item The working directory \texttt{/work} is meant to be a bind mount, so the
        source and resulting PDF live on the host rather than inside the container.
  \item \texttt{ENTRYPOINT} runs \texttt{latexmk}, so extra arguments (such as a
        different file name) can be passed to \texttt{docker run}.
\end{itemize}

\section{Build and Run}
\begin{listing}[H]
\begin{minted}[fontsize=\small,breaklines,frame=lines]{bash}
# Build the image
docker build -t latex-docker .

# Compile hello.tex; the current directory is mounted at /work
docker run --rm -v "$PWD":/work latex-docker hello.tex

# The PDF is now on the host
ls -l hello.pdf
\end{minted}
\caption{Building the image and compiling the document.}
\end{listing}

A helper script wraps the same steps:

\begin{listing}[H]
\inputminted[fontsize=\small,breaklines,frame=lines]{bash}{latex-docker/compile.sh}
\caption{\texttt{compile.sh}: build-and-compile helper.}
\end{listing}

\section{Compiling This Report}
Because the Dockerfile uses \texttt{latexmk} with \texttt{-shell-escape}, it can
also build documents that use the \texttt{minted} package, such as this report.
\texttt{minted} additionally needs Python's Pygments, which the Dockerfile
installs. To compile the whole report:

\begin{listing}[H]
\begin{minted}[fontsize=\small,breaklines,frame=lines]{bash}
docker run --rm -v "$PWD":/work latex-docker devopsAssignment.tex
\end{minted}
\caption{Compiling a minted-based document.}
\end{listing}

\section{Retrospective}
Running the compiler in a container gives a reproducible build: the same TeX
Live version produces the same PDF regardless of the host operating system.
The tradeoff is image size and build time on first use, which is why the image
installs only the packages the document needs.
